% !TEX program = xelatex
\documentclass{ctexart}
\ctexset{section/format=\Large\bfseries}
\usepackage[svgnames]{xcolor}
\usepackage[a4paper, margin = 2.33cm, centering, includefoot]{geometry}
%\linespread{1.1}
\usepackage{titletoc}
\titlecontents{section}[1.8em]{}{\sffamily\contentslabel{1.5em}}{}{\,\titlerule*[0.5pc]{$\cdot$}\contentspage}[\vspace{1.1ex}]
%\AddToHook{cmd/section/before}{\clearpage}
%\renewcommand{\thesection}{\S\arabic{section}}
\usepackage[unicode=true, bookmarksnumbered=true, colorlinks,
	citecolor=ForestGreen,       % color of links to bibliography
	linkcolor=Magenta,		     % color of internal links
	urlcolor=NavyBlue	     	 % color of external links
]{hyperref}
\usepackage{graphicx,subfig}
\usepackage{enumerate}
\usepackage{amsthm,amssymb,amsmath,mathrsfs}
\usepackage{bm,bbm,extarrows,cancel,manfnt,marvosym}

\newtheorem{theorem}{Theorem}
\newtheorem{definition}{Definition}
\newtheorem*{note}{Note}
\newtheorem*{remark}{Remark}
\newtheorem{example}{Example}
\newtheorem{lemma}{Lemma}
\newtheorem{corollary}{Corollary}
\newtheorem{proposition}{Proposition}

\newcommand{\textred}[1]{\textcolor{red}{#1}}


\title{\textsf{Data-Driven}}

\begin{document}
\pagenumbering{roman}
\setcounter{page}{0}

% \pdfbookmark[1]{随机分析}{title}
% \pdfbookmark[1]{随机分析}{title} 命令会在 PDF 文档中创建一个名为“随机分析”的次级书签。当你打开生成的 PDF 文件并查看书签面板时，你会看到一个名为“随机分析”的书签。如果有其他的书签，它们会按照层级关系显示在一起。
\maketitle

\vspace{-2ex}\par\hspace{1.5em}
概率统计、计算数学和应用数学组的研究课题与统计、流体力学、材料科学、群体运动、医学成像、图像处理、机器学习等密切相关;
数学基础和研究方法包括概率、偏微分方程、调和分析、几何、科学计算、数值分析、优化算法等

\pdfbookmark[1]{封面}{cover}
\vspace{2.333ex}
\vspace{-2.333ex}
\clearpage

\pdfbookmark[1]{\contentsname}{toc}
\tableofcontents

\vspace{-2.333ex}

\clearpage
\pagenumbering{arabic}
\setcounter{page}{1}

%\addtocounter{section}
% 如果你想在当前章节编号的基础上增加一个额外的章节，可以使用如下命令：
% \addtocounter{section}{1}
% 如果当前章节是第 3 章，这个命令将把章节编号调整为第 4 章。这个命令不会影响实际的章节内容，只是改变章节计数器的值。

\section{数学与数据科学：小波的故事}

小波理论和压缩感知是应用数学领域的两个关键研究方向，也是图像科学
中不可或缺的数学工具。小波，作为一类特殊函数，能够高效地生成函数空间
的基或框架，从而有效地表征函数。而压缩感知则作为一种信号采集与重建的
方法论，旨在在不降低信号重建质量的前提下，尽可能减少所需的采样数据。

数据科学，作为一门集统计学、计算机科学和信息 理论于一体的跨学科领域，专注于从大规模数据集中挖掘知识。


{\color{red} 有data -> 数据科学 //
No data, just theorem ->数学 }

迈耶尔（Yves Meyer）和马拉（Stéphane Mallat）提出的多分辨率分析
（multiresolution analysis, MRA）理论，对小波分析领域产生了革命性的影响。
多分辨率分析不仅桥接了离散和连续小波变换，还极大地简化了小波的构建过程。多分辨率分析理论不仅让迈耶尔构造出了在频域具有紧支集的band limited小波，也为道贝希斯（Ingrid Daubechies）的正交小波的发展铺平了道路。道贝希斯小波以其出色的平滑性和紧支集特性，在数据处理中展示了显著优势。

小波分析与数据处理的密切关系主要体现在信号和图像处理方面。尤其是多分辨率分析理论中提出的多尺度分析，使小波变换能够更有效地适应不同尺度的数据特征。多分辨率分析构造的小波对应着滤波器组，可以通过多重卷积运算来实现快速小波变换算法（Fast Wavelet Transform, FWT），为小波分析在信号和图像处理中的应用创建了独特的优势。小波分析不仅为信号提供了多尺度表达，还能够利用其稀疏逼近性质从而高效地进行信号和图像的压缩与重建。小波的这些独特的数学特性，使它成为了数据处理的一个强大的工具。

然而，随着时间的推进，傅里叶分析在处理非光滑或突变信号时的局限性逐渐暴露。这种局限性激发了研究人员对替代信号分析方法的探索，尤其是在 20 世纪 70 年代和 80 年代，小波分析开始获得日益增长的关注，特别是对小波在时域和频域中提供局部信息的能力。

调和分析——一种将复杂数学对象分
 解成简单波状成分的方法
 
\subsection{多分辨率分析}
多分辨率分析是小波的理论和应用的基石
。多分辨率分析的核心在于简
化了具备理想特性小波的构造过程，使得小波变得更加易于构造和应用。这一
概念不仅成为小波理论中的一个重要组成部分，也使得基于多分辨率分析构造
的小波系统自然地具备了快速的正变换和反变换算法。这些算法是由小波函数
对应的滤波器组组成的，通过多层次的线性滤波操作实现。多分辨率分析的这
一特性不仅让小波构造变得更加容易，同时也自然地建立起离散和连续之间的
桥梁，因此，通过多分辨率分析构造出的小波自然具有快速小波变换。
所谓的多分辨率分析，就是一系列嵌套的子空间，他们把完整的函数空间
按照逐步升高的分辨率来划分，使得构造小波基变得更加容易。数学上来说，
多分辨率分析由子空间 \(V_n \subset H, n \in \mathbf{Z} \)组成，这里\(H\)是某个希尔伯特空间，这些子空间还要满足下述条件:
\[ V_n \subset V_{n+1}, \cap_n V_n = \{0 \}, \overline{\cup_n V_n}=H, f \in V_0 \iff f(2^n \cdot ) \in V_n  \]
且 V0 由一个函数（可细分函数，refinable function）和它的平移所张成，且这些
平移之间是相互正交的。正交小波函数就是 V0 在 V1 的正交补空间的生成元。


\section{黎曼重排定理 Rearranging the Alternating Harmonic Series}

加法交换律：  a + b = b + a 
more generally, 

\begin{align*}
\text{sum}\{a_1, a_2, \dots, a_n\} &= a_1 + a_2 + \dots + a_n \\
&= a_n + a_1 + a_2 + \dots + a_{n-1} \\
&= a_2 + a_1 + a_n + a_{n-1} + \dots + a_3 \\
&= \text{etc.}
\end{align*}


也就是说，交换律说明在有限和情况下重新排列和式不改变总量；或许令人感到吃惊，in an infinite sum 重新排列和式能够改变总量

fact:

\[1 - \frac{1}{2} + \frac{1}{3} - \frac{1}{4} + \dots = ln2 \]

\begin{align*}
 & 1 - \frac{1}{2} + \frac{1}{3} - \frac{1}{4} + \dots  \\
&= (1 - \frac{1}{2}) - \frac{1}{4} + (\frac{1}{3} - \frac{1}{6}) - \frac{1}{8} \\
&= \frac{1}{2} - \frac{1}{4} + \frac{1}{6} - \frac{1}{8} + \dots \\
&= \frac{1}{2} \left(1 - \frac{1}{2} + \frac{1}{3} - \frac{1}{4} + \dots\right) \\
&= \frac{1}{2} \ln 2
\end{align*}

\begin{theorem}[Riemann Theorem]

A conditional convergent series(of real numbers) can be arranged to sum to any real number.

\end{theorem}

\begin{definition}[rearrangement 重排]
We say the series \( \sum b_m \) is a rearrangement of the series \( \sum a_n \) if each
term of the series \( \sum b_m \) occurs as a term of the series \( \sum a_n \) (and the same
number of times) and vice versa.
\end{definition}

我们说级数 \( \sum b_m \) 是级数 \( \sum a_n \) 的重排，如果 \( \sum b_m \) 中的每一项都出现在 \( \sum a_n \) 中（并且出现的次数相同），反之亦然

Outline of proof of Riemann’s Theorem.

\begin{itemize}
\item Since the given series converges, the terms of the series are small.
\item Since the given series is conditionally convergent, the series of positive termsdiverges and the series of negative terms diverges. 
\item Given S, we want to form a rearrangement of the series that sums to S.
\item To begin, select enough of the positive terms (choosing largest first) until the partial sum is larger than S.(This is possible!) 
\item For the next terms of the series, choose enough of the negative terms (choosing the most negative terms first) to make the partial sum less than S. (This is possible!) 
\item Continue, choosing the positive and negatives terms in order of decreasing size so that the partial sums swing to more than S, then less, then more, etc. 
\item Since the terms of the series are small, the oscillations get smaller and smaller so that the sequence of partial sums of the rearranged series converges to S. 
\end{itemize}


\begin{definition}[simple rearrangement 简单重排]
A simple rearrangement of a series is a rearrangement of the series in which the positive terms of the rearranged series occur in the same order as the original series and the negative terms occur in the same order.
\end{definition}

\begin{example}

\[ 1- \frac{1}{2} - \frac{1}{4} + \frac{1}{3} - \frac{1}{6} -\frac{1}{8} + \dots \]是AHS的simple rearrangement,但是
\[ \frac{1}{3} - \frac{1}{2} -\frac{1}{6} + 1 + \frac{1}{5} - \frac{1}{4} + \frac{1}{7} - \frac{1}{12} + \dots \] 不是AHS的simple rearrangement

\end{example}

Goal:Describe the sum of every simple rearrangement of the Alternating Harmonic Series.\\
To do this, we’ll use power series.\\
A power series (centered at 0) is a function of the form \[ f(x)= \sum\limits_{n=0}^{\infty}a_nx^n = a_0 + a_1x + a_2x^2 + a_3 x^3 + \dots  \]

If the series converges for any non-zero x, there is an R > 0 so that the series converges in the open interval −R < x < R. In this interval, the series can be differentiated and integrated term by term and the resulting series also converge in this open interval.

\begin{theorem}[Abel's theorems]
如果 \(\sum a_n\)收敛，并且如果\(f(x) = \sum a_n x^n \),那么 \[ \sum a_n = \lim_{x \to 1^-} f(x) \]
\end{theorem}

\begin{definition}[asymptotic density 渐进密度]

If \(\sum a_n \) is a simple rearrangement of the Alternating Harmonic Series, let \(p_k\) be the number of positive terms in the first k terms, \{ \(a_1, a_2, a_3, \dots , a_k \) \}. The asymptotic density, \( \alpha \), of the positive terms in the rearrangement is \( \alpha = \lim_{k \to \infty} \frac{p_k}{k} \) if the limit exists
\end{definition}

In the rearrangement 

\[
1 - \frac{1}{2} - \frac{1}{4} + \frac{1}{3} - \frac{1}{6} - \frac{1}{8} + \frac{1}{5} - \dots
\]

\[
\begin{array}{c|c|c}
k & \{a_1, \dots, a_k\} & p_k \\
\hline
1 & \{1\} & 1 \\
2 & \{1, -\frac{1}{2}\} & 1 \\
3 & \{1, -\frac{1}{2}, -\frac{1}{4}\} & 1 \\
4 & \{1, -\frac{1}{2}, -\frac{1}{4}, \frac{1}{3}\} & 2 \\
5 & \{1, -\frac{1}{2}, -\frac{1}{4}, \frac{1}{3}, -\frac{1}{6}\} & 2 \\
  & \text{Etc.}  & \\
\end{array}
\]

so the asymptotic density is \(\frac{1}{3}\).

\begin{theorem}
对于交错调和级数的一个简单重排，当且仅当重排中的正项的渐近密度 \(\alpha\) 存在时，该级数收敛到一个扩展实数(an extended real number)

此外，具有渐近密度 \(\alpha\) 的重排的和为：

\[
\ln 2 + \frac{1}{2} \ln \left(\frac{\alpha}{1 - \alpha}\right)
\]
\end{theorem}

Similarly, in the rearrangement

\[
1 + \frac{1}{3} - \frac{1}{2} + \frac{1}{5} + \frac{1}{7} - \frac{1}{4} + \frac{1}{9} + \frac{1}{11} - \frac{1}{6} + \dots
\]

\[
p_1 = 1, \quad p_2 = 2, \quad p_3 = 2, \quad p_4 = 3, \quad p_5 = 4, \quad p_6 = 4, \dots
\]

so the asymptotic density is \(\alpha = \lim_{k \to \infty} \frac{p_k}{k} = \frac{2}{3}\)

and the sum of the series is

\[
\ln 2 + \frac{1}{2} \ln \left(\frac{\frac{2}{3}}{1 - \frac{2}{3}}\right) = \ln 2 + \frac{1}{2} \ln 2 = \frac{3}{2} \ln 2
\]

\begin{proof}

假设 \(\sum_{n=1}^{\infty} a_n\) 是一个简单的 AHS（调和级数）重新排列。

令 \(p_k\) 表示前 \(k\) 项中正项的数量。那么在前 \(k\) 项中负项的数量为 \(q_k = k - p_k\)。

由于给定的级数是 AHS 的一个简单重新排列，我们有：

\[
\sum_{n=1}^{k} a_n = \sum_{j=1}^{p_k} \frac{1}{2j-1} - \sum_{j=1}^{q_k} \frac{1}{2j}
\]

对于每一个正整数 \(m\)，令

\[
E_m = \sum_{n=1}^{m} \frac{1}{n} - \ln m
\]

序列 \(E_1, E_2, E_3, \dots\) 是一个正的递减序列，其极限 \(\gamma \approx 0.5772\) 称为欧拉常数。

由于

\[
\sum_{n=1}^{m} \frac{1}{n} = \ln m + E_m
\]

我们有：

\[
\sum_{j=1}^{q_k} \frac{1}{2j} = \frac{1}{2} \sum_{j=1}^{q_k} \frac{1}{j} = \frac{1}{2} \ln q_k + \frac{1}{2} E_{q_k}
\]

并且

\[
\sum_{j=1}^{p_k} \frac{1}{2j-1} = \sum_{\ell=1}^{2p_k} \frac{1}{\ell} - \sum_{\ell=1}^{p_k} \frac{1}{2\ell} = \ln(2p_k) + E_{2p_k} - \frac{1}{2} \ln p_k - \frac{1}{2} E_{p_k}
\]

因此，

\[
\sum_{n=1}^{k} a_n = \ln(2p_k) + E_{2p_k} - \frac{1}{2} \ln p_k - \frac{1}{2} E_{p_k} - \frac{1}{2} \ln q_k - \frac{1}{2} E_{q_k}
\]

经过简化，我们得到：

\[
\sum_{n=1}^{k} a_n = \ln 2 + \frac{1}{2} \ln \left(\frac{p_k}{q_k}\right) + E_{2p_k} - \frac{1}{2} E_{p_k} - \frac{1}{2} E_{q_k}
\]

这是重新排列后级数的第 \(k\) 部分和的精确公式。

要找到重新排列后级数的和，我们需要取部分和的极限。

注意：

\[
\frac{p_k}{q_k} = \frac{p_k}{k - p_k} = \frac{p_k}{k} \cdot \frac{1}{1 - \frac{p_k}{k}}
\]

因此，如果 \(\alpha\) 存在，

\[
\lim_{k \to \infty} \frac{p_k}{q_k} = \frac{\alpha}{1 - \alpha}
\]

并且当 \(\alpha\) 不存在时，\(\frac{p_k}{q_k}\) 没有极限。

还注意到：

\[
\lim_{k \to \infty} E_{2p_k} = \lim_{k \to \infty} E_{p_k} = \lim_{k \to \infty} E_{q_k} = \gamma
\]

将这些结果结合在一起，我们得到：

\[
\sum_{n=1}^{\infty} a_n = \lim_{k \to \infty} \sum_{n=1}^{k} a_n = \lim_{k \to \infty} \left(\ln 2 + \frac{1}{2} \ln \left(\frac{p_k}{q_k}\right) + E_{2p_k} - \frac{1}{2} E_{p_k} - \frac{1}{2} E_{q_k}\right)
\]

最终，我们得到：

\[
\sum_{n=1}^{\infty} a_n = \ln 2 + \frac{1}{2} \ln \left(\frac{\alpha}{1 - \alpha}\right) + \gamma - \frac{1}{2}\gamma - \frac{1}{2}\gamma = \ln 2 + \frac{1}{2} \ln \left(\frac{\alpha}{1 - \alpha}\right)
\]

\end{proof}

\begin{remark}[英文的证明]

Suppose \(\sum_{n=1}^{\infty} a_n\) is a simple rearrangement of the Alternating Harmonic Series (AHS).

Let \(p_k\) represent the number of positive terms among the first \(k\) terms \(\{a_1, a_2, a_3, \dots, a_k\}\). Consequently, the number of negative terms in the first \(k\) terms is \(q_k = k - p_k\).

Since the series is a simple rearrangement of the AHS, we have:

\[
\sum_{n=1}^{k} a_n = \sum_{j=1}^{p_k} \frac{1}{2j-1} - \sum_{j=1}^{q_k} \frac{1}{2j}
\]

For each positive integer \(m\), define:

\[
E_m = \sum_{n=1}^{m} \frac{1}{n} - \ln m
\]

The sequence \(E_1, E_2, E_3, \dots\) is a decreasing sequence of positive numbers, with the limit \(\gamma \approx 0.5772\), known as Euler's constant.

Since

\[
\sum_{n=1}^{m} \frac{1}{n} = \ln m + E_m
\]

we can express the sum of the negative terms as:

\[
\sum_{j=1}^{q_k} \frac{1}{2j} = \frac{1}{2} \sum_{j=1}^{q_k} \frac{1}{j} = \frac{1}{2} \ln q_k + \frac{1}{2} E_{q_k}
\]

and the sum of the positive terms as:

\[
\sum_{j=1}^{p_k} \frac{1}{2j-1} = \sum_{\ell=1}^{2p_k} \frac{1}{\ell} - \sum_{\ell=1}^{p_k} \frac{1}{2\ell} = \ln(2p_k) + E_{2p_k} - \frac{1}{2} \ln p_k - \frac{1}{2} E_{p_k}
\]

Thus, we have:

\[
\sum_{n=1}^{k} a_n = \ln(2p_k) + E_{2p_k} - \frac{1}{2} \ln p_k - \frac{1}{2} E_{p_k} - \frac{1}{2} \ln q_k - \frac{1}{2} E_{q_k}
\]

Simplifying this expression, we get:

\[
\sum_{n=1}^{k} a_n = \ln 2 + \frac{1}{2} \ln \left(\frac{p_k}{q_k}\right) + E_{2p_k} - \frac{1}{2} E_{p_k} - \frac{1}{2} E_{q_k}
\]

This gives the exact formula for the \(k\)-th partial sum of the rearranged series.

To find the sum of the rearranged series, we need to take the limit of these partial sums.

Note that:

\[
\frac{p_k}{q_k} = \frac{p_k}{k - p_k} = \frac{p_k}{k} \cdot \frac{1}{1 - \frac{p_k}{k}}
\]

Therefore, if \(\alpha\) exists, we have:

\[
\lim_{k \to \infty} \frac{p_k}{q_k} = \frac{\alpha}{1 - \alpha}
\]

If \(\alpha\) does not exist, \(\frac{p_k}{q_k}\) does not have a limit.

Additionally, note that:

\[
\lim_{k \to \infty} E_{2p_k} = \lim_{k \to \infty} E_{p_k} = \lim_{k \to \infty} E_{q_k} = \gamma
\]

Combining all these results, we obtain:

\[
\sum_{n=1}^{\infty} a_n = \lim_{k \to \infty} \sum_{n=1}^{k} a_n = \lim_{k \to \infty} \left(\ln 2 + \frac{1}{2} \ln \left(\frac{p_k}{q_k}\right) + E_{2p_k} - \frac{1}{2} E_{p_k} - \frac{1}{2} E_{q_k}\right)
\]

Finally, we get:

\[
\sum_{n=1}^{\infty} a_n = \ln 2 + \frac{1}{2} \ln \left(\frac{\alpha}{1 - \alpha}\right) + \gamma - \frac{1}{2}\gamma - \frac{1}{2}\gamma = \ln 2 + \frac{1}{2} \ln \left(\frac{\alpha}{1 - \alpha}\right)
\]



\end{remark}


黎曼关于条件收敛级数重新排列的定理是非构造性的。相比之下，我们已经看到，AHS 的简单重新排列当且仅当正项的渐近密度存在时才会收敛，并且当渐近密度存在时，我们有一个公式来求级数的和。特别地，AHS 的简单重新排列的和仅依赖于正项的渐近密度。

\subsection{压缩感知(compressed sensing)}

在传统模式下，先完成采样再进行数据压缩的做法效率并不高。这激发了一个新的思考 ：能否在采样阶段就考虑到图像的稀疏特性，从而在数据采集时就直接实现压缩呢？压缩感知的概念应运而生，其核心在于在不损失数据重建质量的前提下，减少所需的采样次数。然而，如何在具体操作上实现这一目标，并确保有理论基础和高效算法的支撑，仍是需要解决的关键科学问题。

\section{SVD降维}

当 \( n \ge m \) 时，矩阵 \( \Sigma \) 至多在对角线上有 \( m \) 个非零元素，可以写成如下形式：

\[ \Sigma = \begin{pmatrix}
\hat{\Sigma} \\
0
\end{pmatrix} \]

因此，可以使用经济型SVD（economy SVD）准确地表示 \( X \)：

\[ X = U \Sigma V^* = \begin{pmatrix}
\hat{U} & \hat{U}_\perp
\end{pmatrix} \begin{pmatrix}
\hat{\Sigma} \\
0
\end{pmatrix} V^* = \hat{U} \hat{\Sigma} V^* \]

这里，\( \hat{U} \) 是一个 \( n \times m \) 的矩阵，包含了 \( U \) 的前 \( m \) 列，\( \hat{\Sigma} \) 是一个 \( m \times m \) 的对角矩阵，包含了 \( \Sigma \) 的前 \( m \) 行和列，\( \hat{U}_\perp \) 是一个 \( n \times (n - m) \) 的矩阵，包含了 \( U \) 的剩余列。


\subsection{SVD证明}

\href{https://blog.csdn.net/v20000727/article/details/135682784}{奇异值分解(SVD)【详细推导证明】}



\subsection{SVD计算过程}

\textbf{Step 1}
根据spectral theorem

The spectral theorem extends to a more general class of matrices. Let \( A \) be an operator on a finite-dimensional inner product space. \( A \) is said to be normal if \( A^* A = A A^* \).

One can show that \( A \) is normal if and only if it is unitarily diagonalizable using the Schur decomposition. That is, any matrix can be written as \( A = U T U^* \), where \( U \) is unitary and \( T \) is upper triangular. If \( A \) is normal, then one sees that \( T T^* = T^* T \). Therefore, \( T \) must be diagonal since a normal upper triangular matrix is diagonal. The converse is obvious.

In other words（换句话说）, \( A \) is normal if and only if there exists a unitary matrix \( U \) such that

\[ A = U D U^*, \]

where \( D \) is a diagonal matrix. Then, the entries of the diagonal of \( D \) are the eigenvalues of \( A \). The column vectors of \( U \) are the eigenvectors of \( A \), and they are orthonormal. Unlike the Hermitian case, the entries of \( D \) need not be real.

\textbf{Step 2}
此外，\( AA^*\)和\( A^*A\)有相同的特征值和特征向量，可以利用线性方程组来帮助证明.

\textbf{Step 3}

The singular value decomposition (SVD) is a versatile tool that can be applied to any \( m \times n \) matrix, unlike eigenvalue decomposition, which is limited to square diagonalizable matrices. Despite this difference, SVD and eigenvalue decomposition are related in the following ways:

If \( \mathbf{M} \) has an SVD \( \mathbf{M} = \mathbf{U} \mathbf{\Sigma} \mathbf{V}^* \), the following relations hold:

\[ 
\mathbf{M}^* \mathbf{M} = \mathbf{V} \mathbf{\Sigma}^* \mathbf{U}^* \mathbf{U} \mathbf{\Sigma} \mathbf{V}^* = \mathbf{V} (\mathbf{\Sigma}^* \mathbf{\Sigma}) \mathbf{V}^*, 
\]

\[ 
\mathbf{M} \mathbf{M}^* = \mathbf{U} \mathbf{\Sigma} \mathbf{V}^* \mathbf{V} \mathbf{\Sigma}^* \mathbf{U}^* = \mathbf{U} (\mathbf{\Sigma} \mathbf{\Sigma}^*) \mathbf{U}^*. 
\]

The right-hand sides of these relations describe the eigenvalue decompositions of the left-hand sides. Consequently:

1. The columns of \( \mathbf{V} \) (right-singular vectors) are eigenvectors of \( \mathbf{M}^* \mathbf{M} \).

2. The columns of \( \mathbf{U} \) (left-singular vectors) are eigenvectors of \( \mathbf{M} \mathbf{M}^* \).

3. The non-zero elements of \( \mathbf{\Sigma} \) (non-zero singular values) are the square roots of the non-zero eigenvalues of \( \mathbf{M}^* \mathbf{M} \) or \( \mathbf{M} \mathbf{M}^* \).

\begin{example}
    参考 \quad \href{https://www.cnblogs.com/marsggbo/p/10155801.html}{奇异值分解(SVD)计算过程示例}
\end{example}

\section{FFT和小波}

自然地，在数据具有无限精细分辨率的极限情况下，离散和连续的形式应当一致。傅里叶级数和傅里叶变换与无限维函数空间（或称希尔伯特空间）的几何紧密相关，希尔伯特空间将向量空间的概念推广到包含无限多自由度的函数。(Naturally, the discrete and continuous formulations should match in the limit of data with infinitely fine resolution. The Fourier series and transform are intimately related to the geometry of infinite-dimensional function spaces, or Hilbert spaces, which generalize the notion of vector spaces to include functions with infinitely many degrees of freedom.)

\vspace{6ex} %此命令用于上下间距的调整

\begin{remark}[infinitely fine resolution]

“无限精细分辨率”可以理解为数据的分辨率达到无限小的极限。在这种情况下，数据点之间的间隔趋近于零，数据可以被看作是连续的，而不再是离散的。这意味着我们拥有无限多的数据点来精确描述一个函数或信号。在数学上，这种情况与处理连续函数的分析方法（如傅里叶级数和傅里叶变换）相对应。总之，“无限精细分辨率”表示在极限情况下，离散数据变得如此精细，以至于它们可以被认为是连续的。

\end{remark}

\vspace{6ex} %此命令用于上下间距的调整

\textcolor{red}{特别地，如果我们将函数 \( f(x) \) 和 \( g(x) \) 离散化为数据向量，我们希望随着采样分辨率的提高，向量内积能够收敛于函数内积。(In particular, if we discretize the functions f(x) and g(x) into vectors of data, as in Fig. 2.1, we would like the vector inner product to converge to the function inner product as the sampling resolution is increased.)}

就像在有限维向量空间中一样，内积可以用来将一个函数投影到由一组正交函数定义的新坐标系中。一个函数 \( f \) 的傅里叶级数表示正是将这个函数投影到在区间 \([a,b]\) 上具有整数周期的正弦和余弦函数的正交集合上。


\subsection{Fourier级数}

\begin{lemma}

Bessel不等式 \quad Parsevel 等式（勾股定理）

\end{lemma}

\begin{definition}[piecewise smooth]

\end{definition}

\begin{theorem}[Fourier级数收敛定理]

周期；分段光滑

\end{theorem}

截断的傅里叶级数在阶跃函数的尖锐拐角处会受到振铃振荡的困扰，这种现象被称为 Gibbs 现象.(The truncated Fourier series is plagued by \textred{ringing oscillations}, known as Gibbs phenomena, around the sharp corners of the step function.)

\begin{remark}{Gibbs现象}


\textbf{Gibbs现象} 是指在傅里叶级数的截断（或截断傅里叶级数）近似中，出现的振铃效应或振荡现象。这种现象常见于对具有尖锐变化（如阶跃函数）的函数进行傅里叶级数近似时。具体表现为，在函数的急剧变化（例如阶跃的边界）处，傅里叶级数的近似会出现振荡，振幅在尖锐的边界附近达到一个明显的峰值，这种现象被称为 \textbf{Gibbs现象}。

\textbf{Gibbs现象的特点}
\begin{enumerate}
    \item \textbf{振荡行为}：在函数的跳跃点（例如阶跃函数的边界）附近，傅里叶级数的近似表现为过度的振荡。这些振荡不会随着更多的傅里叶级数项的增加而消失，只会在振幅上有所降低，但不会完全消失。
    \item \textbf{振幅}：尽管振荡会随着傅里叶级数项数目的增加而减小，振幅不会完全收敛到零。通常在尖锐边界附近，振幅约为真实跳跃幅度的 9\% 左右。
    \item \textbf{影响因素}：Gibbs现象的强度取决于原函数的尖锐性和傅里叶级数的截断程度。对于阶跃函数或其他具有急剧变化的函数，Gibbs现象尤其明显。
\end{enumerate}

\textbf{例子：阶跃函数}
考虑阶跃函数（Heaviside阶跃函数）作为例子。它在某一点（通常是 $x=0$）有一个不连续的跳跃。傅里叶级数试图通过有限的正弦和余弦函数的线性组合来近似阶跃函数，这会导致在阶跃点附近出现显著的振荡，这就是Gibbs现象。

\textbf{解决方法}
\begin{itemize}
    \item \textbf{平滑处理}：可以使用平滑处理技术来减轻Gibbs现象，例如使用窗函数来减少高频成分。
    \item \textbf{正则化}：在某些应用中，可以通过正则化方法来控制振荡。
    \item \textbf{其他变换}：使用其他变换方法，如小波变换，可能会减少Gibbs现象的影响。
\end{itemize}

总的来说，Gibbs现象是傅里叶级数截断的一个固有问题，了解和处理这些振荡对于信号处理和数值分析等领域非常重要。


\end{remark}


\subsection{离散傅里叶变换}

尽管我们总是使用 FFT 进行计算，但从离散傅里叶变换（DFT）的最简单形式开始是有益的。离散傅里叶变换定义为：

\[
\hat{f}_k = \sum_{j=0}^{n-1} f_j e^{-i2\pi jk/n} \quad (2.26)
\]

而离散傅里叶变换的逆变换（iDFT）定义为：

\[
f_k = \frac{1}{n} \sum_{j=0}^{n-1} \hat{f}_j e^{i2\pi jk/n} \quad (2.27)
\]

\textbf{\color{red}因此，DFT 是一个线性算子（即矩阵），将数据点 \( f \) 映射到频域 \( \hat{f} \)}：

\[
\{f_1, f_2, \cdots, f_n\} \xrightarrow{\text{DFT}} \{ \hat{f}_1, \hat{f}_2, \cdots, \hat{f}_n \} \quad (2.28)
\]

对于给定数量的点 \( n \)，DFT 使用正弦和余弦函数表示数据，这些函数的频率是基本频率的整数倍， \( \omega_n = e^{-2\pi i/n} \)。DFT 可以通过矩阵乘法计算：

\[
\begin{bmatrix}
\hat{f}_1 \\
\hat{f}_2 \\
\hat{f}_3 \\
\vdots \\
\hat{f}_n
\end{bmatrix}
=
\begin{bmatrix}
1 & 1 & \cdots & 1 \\
1 & \omega_n & \omega_n^2 & \cdots & \omega_n^{n-1} \\
1 & \omega_n^2 & \omega_n^4 & \cdots & \omega_n^{2(n-1)} \\
\vdots & \vdots & \vdots & \ddots & \vdots \\
1 & \omega_n^{n-1} & \omega_n^{2(n-1)} & \cdots & \omega_n^{(n-1)(n-1)}
\end{bmatrix}
\begin{bmatrix}
f_1 \\
f_2 \\
f_3 \\
\vdots \\
f_n
\end{bmatrix}
\quad (2.29)
\]

输出向量 \( \hat{f} \) 包含输入向量 \( f \) 的傅里叶系数，DFT 矩阵 \( F \) 是一个单位 Vandermonde 矩阵。矩阵 \( F \) 是复值的，因此输出 \( \hat{f} \) 具有幅度和相位，这两者都有有用的物理解释。DFT 矩阵 \( F \) 的实部在图 2.8 中显示 \( n = 256 \)。代码 2.2 生成并绘制了这个矩阵。从图像中可以看出，\( F \) 存在分层且高度对称的多尺度结构。每一行和每一列都是频率逐渐增加的余弦函数。

\subsection{快速傅里叶变换}

如前所述，通过 DFT 矩阵 \( F \) 进行乘法涉及 \( O(n^2) \) 操作。快速傅里叶变换（FFT）缩放为 \( O(n \log(n)) \)，使得它可以广泛应用于各种场景，包括 MP3 和 JPG 格式的音频和图像压缩、视频流、卫星通信以及蜂窝网络，仅仅是这些众多应用的一部分。

例如，音频通常以 44.1 kHz 采样，即每秒 44,100 次采样。对于 10 秒的音频，向量 \( f \) 将具有维度 \( n = 4.41 \times 10^5 \)。使用矩阵乘法计算 DFT 涉及大约 \( 2 \times 10^{11} \) 或 2000 亿次乘法。相比之下，FFT 需要大约 \( 6 \times 10^6 \) 次乘法，这意味着加速因子超过 30,000。因此，FFT 已经成为 DFT 的同义词，FFT 库被内置于几乎所有执行数字信号处理的设备和操作系统中。

要看到 FFT 的巨大好处，可以考虑音频信号的传输、存储和解码。我们将看到许多信号在傅里叶变换域中高度可压缩，这意味着大多数 \( \hat{f} \) 的系数都很小，可以被丢弃。这使得压缩信号的存储和传输更加高效，因为只需要传输非零的傅里叶系数。然而，随后必须通过计算 FFT 和逆 FFT（iFFT）来快速编码和解码压缩的傅里叶信号。

FFT 的基本思想是，当数据点 \( n \) 是 2 的幂时，DFT 可以更高效地实现。例如，考虑 \( n = 1024 = 2^{10} \)。在这种情况下，DFT 矩阵 \( F_{1024} \) 可以写成：

\[
\hat{f} = F_{1024} f =
\begin{bmatrix}
I_{512} & -D_{512} \\
I_{512} & -D_{512}
\end{bmatrix}
\begin{bmatrix}
F_{512} & 0 \\
0 & F_{512}
\end{bmatrix}
\begin{bmatrix}
f_{\text{even}} \\
f_{\text{odd}}
\end{bmatrix}
\quad (2.30)
\]

其中 \( f_{\text{even}} \) 是 \( f \) 的偶数索引元素，\( f_{\text{odd}} \) 是奇数索引元素，\( I_{512} \) 是 \( 512 \times 512 \) 单位矩阵，\( D_{512} \) 定义为：

\[
D_{512} =
\begin{bmatrix}
1 & 0 & \cdots & 0 \\
0 & \omega & 0 & \cdots & 0 \\
0 & 0 & \omega^2 & \cdots & 0 \\
\vdots & \vdots & \vdots & \ddots & \vdots \\
0 & 0 & \cdots & 0 & \omega^{511}
\end{bmatrix}
\quad (2.31)
\]

这个表达式可以通过仔细计算和重组 (2.26) 和 (2.29) 中的项得到。如果 \( n = 2^p \)，这个过程可以重复，\( F_{512} \) 可以表示为 \( F_{256} \)，然后可以表示为 \( F_{128} \rightarrow F_{64} \rightarrow F_{32} \rightarrow \cdots \)。如果 \( n \neq 2^p \)，可以用零填充向量直到它是 2 的幂。FFT 然后涉及对 \( f \) 的子向量的偶数和奇数索引的高效交错，以及多个较小的 2 × 2 DFT 计算。


\subsection{快速傅里叶变换中的分解与合并}

分解和合并是快速傅里叶变换（FFT）降低计算复杂度的核心思想。

\begin{itemize}
    \item \textbf{递归分解}：每次分解将序列长度减少一半，直到达到最小的 DFT（通常是长度为 2 的 DFT）。这种递归分解的深度是 \( \log_2 n \)。每层分解操作的数量是 \( n \)，总的计算复杂度为 \( O(n \log n) \)。
    
    \item \textbf{减少重复计算}：在每一层分解和合并过程中，避免了对相同计算的重复执行。通过合并步骤，利用了之前计算的部分结果，减少了总体的计算量。
\end{itemize}

\subsubsection{具体示例}

假设我们有一个长度为 8 的序列 \( f \)，使用 FFT 进行处理：

\begin{itemize}
    \item \textbf{第一次分解}：
    \begin{itemize}
        \item 偶数索引： \( [f_0, f_2, f_4, f_6] \)
        \item 奇数索引： \( [f_1, f_3, f_5, f_7] \)
    \end{itemize}

    \item \textbf{第二次分解}（分别对每个长度为 4 的序列进行分解）：
    \begin{itemize}
        \item 偶数索引的偶数部分： \( [f_0, f_4] \)
        \item 偶数索引的奇数部分： \( [f_2, f_6] \)
        \item 奇数索引的偶数部分： \( [f_1, f_5] \)
        \item 奇数索引的奇数部分： \( [f_3, f_7] \)
    \end{itemize}

    \item 计算每个长度为 2 的 DFT。

\end{itemize}

\subsection{合并过程}

合并过程是快速傅里叶变换（FFT）中的关键步骤，它通过将较小的 DFT 结果合并成较大的 DFT 结果来减少计算复杂度。具体来说：

\subsubsection{分解后的 DFT}

考虑一个长度为 \( n \) 的序列 \( f \)，我们已经将其分解为两个长度为 \( n/2 \) 的子序列：偶数索引部分 \( f_{\text{even}} \) 和奇数索引部分 \( f_{\text{odd}} \)。这两个子序列的 DFT 结果分别为 \( \hat{f}_{\text{even}} \) 和 \( \hat{f}_{\text{odd}} \)。我们的目标是将这两个结果合并为原始序列长度 \( n \) 的 DFT 结果 \( \hat{f} \)。

\subsubsection{复指数函数的利用}

合并过程中，我们利用复指数函数的对称性和周期性。合并步骤如下：

对于 \( k = 0, 1, \ldots, n/2-1 \)，合并计算公式为：
\[
\hat{f}_k = \hat{f}_{\text{even}, k} + e^{-i2\pi k/n} \hat{f}_{\text{odd}, k}
\]
\[
\hat{f}_{k+n/2} = \hat{f}_{\text{even}, k} - e^{-i2\pi k/n} \hat{f}_{\text{odd}, k}
\]

这里，\( e^{-i2\pi k/n} \) 是复指数因子，用于将偶数部分和奇数部分的结果合并成最终的 DFT 结果。

\subsubsection{合并步骤的详细计算}

\begin{enumerate}
    \item \textbf{计算偶数和奇数部分的 DFT}：首先计算长度为 \( n/2 \) 的偶数部分和奇数部分的 DFT，得到 \( \hat{f}_{\text{even}} \) 和 \( \hat{f}_{\text{odd}} \)。
    
    \item \textbf{应用复指数因子}：计算复指数因子的值，通常为 \( \omega_n^k = e^{-i2\pi k/n} \)，其中 \( \omega_n \) 是基本频率。
    
    \item \textbf{合并结果}：
    \begin{itemize}
        \item 对于 \( k = 0, 1, \ldots, n/2-1 \)：
        \[
        \hat{f}_k = \hat{f}_{\text{even}, k} + \omega_n^k \hat{f}_{\text{odd}, k}
        \]
        \[
        \hat{f}_{k+n/2} = \hat{f}_{\text{even}, k} - \omega_n^k \hat{f}_{\text{odd}, k}
        \]
    \end{itemize}
\end{enumerate}

\subsubsection{例子：长度为 8 的序列}

假设我们有一个长度为 8 的序列 \( f = [f_0, f_1, f_2, f_3, f_4, f_5, f_6, f_7] \)，我们将其分解为两个长度为 4 的子序列：

\begin{itemize}
    \item 偶数索引部分：\( f_{\text{even}} = [f_0, f_2, f_4, f_6] \)
    \item 奇数索引部分：\( f_{\text{odd}} = [f_1, f_3, f_5, f_7] \)
\end{itemize}

然后，分别计算 \( f_{\text{even}} \) 和 \( f_{\text{odd}} \) 的 DFT，得到 \( \hat{f}_{\text{even}} \) 和 \( \hat{f}_{\text{odd}} \)。

合并步骤如下：

\begin{itemize}
    \item \textbf{计算 \( k = 0 \) 的结果}：
    \[
    \hat{f}_0 = \hat{f}_{\text{even}, 0} + \omega_8^0 \hat{f}_{\text{odd}, 0}
    \]
    \[
    \hat{f}_4 = \hat{f}_{\text{even}, 0} - \omega_8^0 \hat{f}_{\text{odd}, 0}
    \]

    \item \textbf{计算 \( k = 1 \) 的结果}：
    \[
    \hat{f}_1 = \hat{f}_{\text{even}, 1} + \omega_8^1 \hat{f}_{\text{odd}, 1}
    \]
    \[
    \hat{f}_5 = \hat{f}_{\text{even}, 1} - \omega_8^1 \hat{f}_{\text{odd}, 1}
    \]

    依此类推，计算 \( k = 2, 3 \) 的结果。
\end{itemize}

\subsubsection{总结}

合并步骤通过将较小的 DFT 结果（偶数和奇数部分的 DFT 结果）应用复指数因子来得到最终的大 DFT 结果。这个过程利用了复数指数函数的对称性和周期性，大大减少了计算的复杂度。这使得 FFT 的计算复杂度从 \( O(n^2) \) 降低到 \( O(n \log n) \)。

\section{稀疏性和压缩感知}

自然数据中观察到的固有结构表明，这些数据在适当的坐标系中可以表示为稀疏的形式。换句话说，如果自然数据用一个精心选择的基底表示，那么仅需要少量参数就可以表征那些处于活跃状态的模式及其比例。所有的数据压缩都依赖于稀疏性，即信号在一个通用的变换基底（如傅里叶基底或小波基底）中通过稀疏系数向量更有效地表示。

最近，数学领域的基础性进展颠覆了这种范式。不再是先收集高维度的测量值然后进行压缩，而是可以直接获取压缩的测量值，并解出与测量值一致的最稀疏的高维信号。

Encoder-Decoder架构

参考\href{https://d2l.ai/chapter_recurrent-modern/encoder-decoder.html}{The Encoder–Decoder Architecture} \qquad 

\href{https://medium.com/@ahmadsabry678/a-perfect-guide-to-understand-encoder-decoders-in-depth-with-visuals-30805c23659b#:~:text=An%20encoder%2Ddecoder%20is%20a,to%20generate%20an%20output%20sequence.}{A Perfect guide to Understand Encoder Decoders in Depth with Visuals}

\href{https://zhuanlan.zhihu.com/p/685837791}{编码器-解码器模型（Encoder-Decoder)}

编码器（Encoder）\qquad 
编码器的作用是接收输入序列，并将其转换成固定长度的上下文向量（context vector）。这个向量是输入序列的一种内部表示，捕获了输入信息的关键特征。在自然语言处理的应用中，输入序列通常是一系列词语或字符。编码器可以是任何类型的深度学习模型，但循环神经网络（RNN）及其变体，如长短期记忆网络（LSTM）和门控循环单元（GRU），因其在处理序列数据方面的优势而被广泛使用。

解码器（Decoder）\qquad 
解码器的目标是将编码器产生的上下文向量转换为输出序列。在开始解码过程时，它首先接收到编码器生成的上下文向量，然后基于这个向量生成输出序列的第一个元素。接下来，它将自己之前的输出作为下一步的输入，逐步生成整个输出序列。解码器也可以是各种类型的深度学习模型，但通常与编码器使用相同类型的模型以保持一致性。

训练过程 \qquad 
在训练Encoder-Decoder模型时，目标是最小化模型预测的输出序列与实际输出序列之间的差异。这通常通过计算损失函数（如交叉熵损失）来实现，并使用反向传播和梯度下降等优化算法进行参数更新。


\section{回归与模型选择}


所有的机器学习都围绕优化展开.
这包括回归和模型选择框架，旨在为数据提供简洁且可解释的模型.

Curve fitting是最基本的回归技术之一，多项式和指数拟合的结果来自于解线性系统Ax = b.

如果模型实现没有被规定好，优化方法被用于选择最好的模型.

这将函数拟合的基础数学(underlying mathematics)改变为如下线性系统的过定或欠定优化问题：

\begin{align}
    \underset{\mathbf{x}}{\arg\min} (|| \mathbf{AX-b}||_2 + \lambda g(\mathbf{x})) \\
    \underset{\mathbf{x}}{\arg\min} g(\mathbf{x}) \text{subject to} 
    ||\mathbf{AX-b}||_2 \leq \epsilon
\end{align}

其中$g(\mathbf{x})$是给定的惩罚项(对于过定系统带有惩罚参数$\lambda$)

\begin{remark}


    这段数学表达式表示两种优化问题的形式：

    1. 带有惩罚项的优化问题：
    \[ \underset{\mathbf{x}}{\arg\min} \left( \| \mathbf{AX} - \mathbf{b} \|_2 + \lambda g(\mathbf{x}) \right) \]
    
    这表示在最小化线性系统残差 \(\| \mathbf{AX} - \mathbf{b} \|_2\) 的同时，也最小化惩罚函数 \(g(\mathbf{x})\)，其中 \(\lambda\) 是惩罚参数，用于控制惩罚项的权重。
    
    2. 带约束的优化问题：
    \[ \underset{\mathbf{x}}{\arg\min} \ g(\mathbf{x}) \ \text{subject to} \ \| \mathbf{AX} - \mathbf{b} \|_2 \leq \epsilon \]
    
    这表示在满足约束条件 \(\| \mathbf{AX} - \mathbf{b} \|_2 \leq \epsilon\) 的前提下，最小化惩罚函数 \(g(\mathbf{x})\)。其中，\(\epsilon\) 是允许的误差范围。
    
    两者的区别在于，第一种方法直接在目标函数中引入惩罚项来控制模型复杂度，而第二种方法则是通过引入约束条件来控制模型误差范围，并在这个约束条件下最小化惩罚函数。
\end{remark}


\vspace{10ex} %此命令用于上下间距的调整

对于方程的未定或者超定线性系统，方程Ax=b的解要么是无解，要么是有无穷多解，一定要选择一个约束或者惩罚，被称之为正则化(regularization)，来得到一个解.

例如，可以在欠定系统中强制使解最小化最小的2范数，从而使得 \(\min g(\mathbf{x}) = \min \| \mathbf{x} \|_2\)。

更一般的，我们考虑回归或者非线性模型，然后整体的数学框架采取更广泛的形式：

\begin{align}
    \underset{\mathbf{x}}{\arg\min} (f(\mathbf{ A,X,b}) + \lambda g(\mathbf{x})) \\
    \underset{\mathbf{x}}{\arg\min} g(\mathbf{x}) \text{subject to} 
    f(\mathbf{ A,X,b}) \leq \epsilon
\end{align}

通常采用梯度下降算法解决.事实上，这个广泛的框架也是深度学习算法的核心.

除了优化策略之外，数据科学的中心关注点是理解一个被提出的模型是否有过拟合或者欠拟合的数据.
Thus cross-validation strategies are
critical for evaluating any proposed model.
交叉验证策略对于评估任何模型都是至关重要的.
A given data set
must be partitioned into a training, validation and withhold set.
训练集、测试集和保留集(withhold set)
根据训练和验证数据建立了一个模型，最后在保留集上进行了测试。对于过拟合，增加模型的复杂度或训练期（迭代）可以提高训练集上的误差，同时导致保留集上的误差增加。

根据训练和验证数据建立了一个模型，最后在保留集上进行了测试。对于过拟合，增加模型的复杂度或训练期（迭代）可以提高训练集上的误差，同时导致保留集上的误差增加。

相反的是，拟合不足(under-fitting)限制了获得一个好的模型的能力.
然而，你的模型是否欠拟合或者模型可以被提升不总是清楚的.

 Importantly, the following mantra holds: if you don’t cross-validate,
you is dumb.

接下来的几章将概述优化和交叉验证在实践中的应用，并强调在应用有意义的约束和结构到 \( g(x) \) 时需要做出的选择，以实现可解释的解。实际上，目标（损失）函数 \( f(\cdot) \) 和正则化 \( g(\cdot) \) 在确定计算上可行的优化策略时同样重要。通常，为了达到优化的真实目标，会选择代理损失和正则化函数。这种选择在很大程度上取决于所考虑的应用领域和数据。


\subsection{经典曲线拟合}

曲线拟合的一个更广泛的数学观点是回归。在本文中，我们将提倡这种观点。与曲线拟合一样，回归试图使用各种统计工具来估计变量之间的关系。具体来说，可以考虑自变量(independent variable) \( \mathbf{X} \)、因变量(dependent variable) \( \mathbf{Y} \) 和一些未知参数 \( \mathbf{\beta} \) 之间的一般关系：

\begin{equation}
    \mathbf{Y} = f(\mathbf{X}, \mathbf{\beta})
\end{equation}

其中，回归函数 \( f(\cdot) \) 通常是预先设定的，参数 \( \beta \) 是通过优化该函数与数据的拟合度来找到的。在接下来的内容中，我们将把曲线拟合作为回归的一种特例。

\vspace{6ex} %此命令用于上下间距的调整

{\color{Red}重要的是，回归和曲线拟合通过优化来发现变量之间的关系。

广义上讲，机器学习围绕回归技术构建，而回归技术本身是基于数据的优化。

因此，从其绝对的数学核心来看，机器学习和数据科学围绕着提出一个优化问题。

当然，优化的成功在很大程度上取决于对要优化的目标函数的定义。}


\subsubsection*{Least-squares 拟合方法}

为了说明回归的概念，我们考虑经典的 least-squares polynomial fitting 去表示数据中的趋势.

这个想法是简单而且直接的:使用一个简单函数去描述趋势，通过最小化被选择的函数和它所拟合的平方和误差.
正如我们所提到的，经典的曲线拟合被表述(formulate)为$\mathbf{Ax=b}$的一个简单解.

考虑一组 \( n \) 个数据点：
\[ (x_1, y_1), (x_2, y_2), (x_3, y_3), \ldots, (x_n, y_n) \ \ (4.5) \]

进一步假设我们希望找到一条最佳拟合线穿过这些点。我们可以用函数：
\[ f(x) = \beta_1 x + \beta_2 \ \ (4.6) \]

来近似这条线，其中常数 \(\beta_1\) 和 \(\beta_2\) 是向量 \(\beta\) 的参数（见公式 (4.4)），它们的选择是为了最小化与拟合相关的一些误差。这条拟合线给出了线性回归模型：
\[ Y = f(A, \beta) = \beta_1 X + \beta_2 \]

因此，该函数提供了一个近似数据的线性模型，每个点的近似误差由下式给出：
\[ f(x_k) = y_k + E_k \ \ (4.7) \]

其中 \( y_k \) 是数据的真实值，\( E_k \) 是拟合与该值的误差。


在用给定函数 \( f(x) \) 进行近似时，可以最小化各种误差度量。在本章中，选择用来计算拟合优度的误差度量或范数将是至关重要的。通常会考虑以下三种标准可能性，它们分别与 \( \ell_2 \)（最小二乘）、\( \ell_1 \) 和 \( \ell_\infty \) 范数相关联。

\[ E_p(f) = \left( \frac{1}{n} \sum_{k=1}^{n} |f(x_k) - y_k|^p \right)^{1/p} \]

对于不同的 \( p \) 值，最佳拟合线将会不同.在大多数情况下，这些差异很小.然而，当数据中存在离群点时，范数的选择可能会产生显著的影响.


\subsubsection*{最小平方直线}
最小二乘拟合线性模型相比其他范数和度量有重要的优势。具体来说，由于误差可以解析计算，优化的成本很低。为了明确展示这一点，考虑将最小二乘拟合准则应用于数据点 \((x_k, y_k)\)，其中 \(k = 1, 2, 3, \ldots, n\)。为了将曲线

\[ f(x) = \beta_1 x + \beta_2 \ \ (4.10) \]

拟合到这些数据，误差 \(E_2\) 通过最小化以下和来找到：

\[ E_2(f) = \sum_{k=1}^{n} |f(x_k) - y_k|^2 = \sum_{k=1}^{n} (\beta_1 x_k + \beta_2 - y_k)^2 \ \ (4.11) \]

最小化这个和需要求导。具体来说，选择常数 \(\beta_1\) 和 \(\beta_2\) 使误差达到最小。因此，我们需要满足：

\[ \frac{\partial E_2}{\partial \beta_1} = 0 \quad \text{和} \quad \frac{\partial E_2}{\partial \beta_2} = 0 \]

请注意，尽管零导数可以表示极小值或极大值，我们知道这一定是误差的最小值，因为没有最大误差，即我们总能选择一条误差更大的线。最小化条件给出：

\[ \frac{\partial E_2}{\partial \beta_1} = 0: \sum_{k=1}^{n} 2(\beta_1 x_k + \beta_2 - y_k)x_k = 0 \ \ (4.12a) \]
\[ \frac{\partial E_2}{\partial \beta_2} = 0: \sum_{k=1}^{n} 2(\beta_1 x_k + \beta_2 - y_k) = 0 \ \ (4.12b) \]

重新排列后，可以得到一个关于 \(\beta_1\) 和 \(\beta_2\) 的 \(2 \times 2\) 线性方程组：

\[ \begin{pmatrix} \sum_{k=1}^{n} x_k^2 & \sum_{k=1}^{n} x_k \\ \sum_{k=1}^{n} x_k & n \end{pmatrix} \begin{pmatrix} \beta_1 \\ \beta_2 \end{pmatrix} = \begin{pmatrix} \sum_{k=1}^{n} x_k y_k \\ \sum_{k=1}^{n} y_k \end{pmatrix} \rightarrow \mathbf{Ax} = \mathbf{b} \ \ (4.13) \]

这个线性方程组可以使用 MATLAB 中的 backslash 命令解决。因此，优化过程是不必要的，因为可以从 \(2 \times 2\) 矩阵中精确计算出解。

这种方法可以很容易地推广到更高次的多项式拟合。特别地，对一组数据进行抛物线拟合需要拟合函数：

\[ f(x) = \beta_1 x^2 + \beta_2 x + \beta_3 \ \ (4.14) \]

现在需要找到三个常数 \(\beta_1, \beta_2, 和 \beta_3\)。通过最小化误差 \(E_2(\beta_1, \beta_2, \beta_3)\) 并求解以下方程组，可以求出这些常数：

\[ \frac{\partial E_2}{\partial \beta_1} = 0 \ \ (4.15a) \]
\[ \frac{\partial E_2}{\partial \beta_2} = 0 \ \ (4.15b) \]
\[ \frac{\partial E_2}{\partial \beta_3} = 0 \ \ (4.15c) \]

实际上，任何 k 次多项式拟合(any polynomial fit of degree k)将产生一个 \((k + 1) \times (k + 1)\) 的线性方程组 \(\mathbf{Ax} = \mathbf{b}\)，可以找到其解。

\subsubsection*{数据线性化}

指数取对数，略.

\subsection{非线性回归与梯度下降}

the root-mean square error 

多项式和指数曲线拟合允许解析上可处理的(analytically tractable)、最佳拟合的最小二乘(best-fit least-square)解。

然而，这种曲线拟合是非常特殊的，为解决更广泛的问题集合，必须采用更通用的数学框架。例如，可能希望将形如 \( f(x) = \beta_1 \cos(\beta_2 x + \beta_3) + \beta_4 \) 的非线性函数拟合到一个数据集上。与求解线性方程组不同，通用的非线性曲线拟合会导致非线性方程组。非线性回归的一般理论假设拟合函数具有如下形式：

\[ f(x) = f(x, \beta) \ \ (4.24) \]


\subsection{回归和Ax=b:超定系统和欠定系统}

Curve fitting result in a optimization problem. -> In many cases, the optimization can be mathematically framed as solving the linear system of equations Ax = b.

许多问题转化为优化问题；

许多优化问题可以表述为解Ax=b问题；

Ax=b超定，约束（方程）比未知量多，通常没有满足线性系统的解，而是找到近似解以最小化给定的误差.

欠定，未知量比约束多，有无穷多解，选择某种约束选择适当而且唯一的解.

本节的目标是hihglight用于优化的两种不同范数 \(l_1\) norm和 \(l_2\) norm，它们用于求解过确定和欠确定系统中的 Ax = b。

    \begin{remark}
        Framework（框架）和Architecture（架构）在软件开发中有不同的含义：

        Architecture（架构）是指软件系统的整体结构和组织方式。它类似于建筑学中的蓝图，规定了系统的功能、接口、数据交换等规范，但不包含具体的程序代码。软件架构通常产生文件、图样、原型和规格，而不是可用的代码。它告诉你系统应该如何构建，但不涉及具体的实现方式.
        
        Framework（框架）是一个已经成形的方法，提供了一套程序代码，告诉你如何使用它。框架通常包含一些基础功能、模块和工具，可以加速开发过程。但它不会告诉你如何实现特定的功能，这是程序员的工作。框架和架构之间的关系是：架构定义了系统的整体结构，而框架则用于实现这个结构的具体细节.
        
        总之，Architecture 是关于系统整体设计的指导方针，而 Framework 是已经成形的方法，帮助你快速开发应用程序。两者相互依赖，共同构建出稳健、高效的软件系统
    \end{remark}


    \vspace{10ex} %此命令用于上下间距的调整
    
    如第3章所述，\(\ell_1\) 范数能够促进解的稀疏性，使得解 \(x\) 的许多载荷为零.( The   1 norm, as already shown in Chapter 3, promotes sparsity so that many of the loadings of the solution x are zero.)

    请注意，添加 \(\ell_1\) 范数会使解变得稀疏，并生成一个以零元素为主的矩阵.(Note that the addition of the  1 norm sparsifies the solution and produces a matrix which is dominated by zero entries.)








    

\subsection{Optimization as the Cornerstone of Regression}

在本章前两节中，拟合函数 \( f(x) \) 已被指定。例如，可能需要生成一个线性拟合，使得 \( f(x) = \beta_1 x + \beta_2 \)。然后，通过之前讨论的回归和优化方法找到这些系数。

接下来，我们的目标是发展技术，允许我们选择一个好的模型去拟合数据，也就是说，我们应该用二次还是三次拟合？

单凭误差指标并不能决定一个好的模型选择，因为选择更多的拟合项会提供更多的参数来降低误差，而不考虑这些额外项是否具有任何意义或可解释性.(The error metric alone does not dictate a good model selection as the more terms that are chosen for fitting, the more parameters are available for lowering the error,  regardless of whether the additional terms have any meaning or interpretability.)

优化策略在从数据中抽取可解释的结果和有意义的模型方面扮演基础性角色.
正如在前面几节所展示的那样，\( \ell_2 \)和\( \ell_1 \) 范数的相互作用在优化结果方面有重要作用.为了进一步阐明优化的作用和可能结果的多样性,
为了进一步说明优化的作用和可能的各种结果，考虑一个简单的例子：数据来自带噪声测量的抛物线
\[ f(x) = x^2 + N(0, \sigma) \]
其中 \( N(0, \sigma) \) 是均值为零、标准差为 \( \sigma \) 的正态分布随机变量。图4.13(a)显示了(4.43)的100次随机测量的一个例子。尽管测量中加入了噪声，但抛物线结构仍然清晰可见。实际上，使用本章第一节中概述的经典最小二乘拟合方法来计算抛物线拟合是很简单的。

{\color{red}目标是发现数据的最佳模型。因此，在实际操作中，我们不知道函数是什么，需要发现它.}

我们可以开始假设一组多项式模型的回归。



%--------------------------------------------

回归结果的可变性对模型选择来说是个问题。

这表明，即使是少量的测量噪声也可能导致对基础模型的显著不同的结论。

接下来，我们将在量化这种可变性的同时，考虑各种解决超定线性系统Ax=b的回归程序。这里强调了五种标准方法：最小二乘回归（pinv），反斜杠操作符（$\backslash$），LASSO（最小绝对收缩和选择算子）（lasso），稳健拟合（robustfit）和岭回归（ridge）.
The variability of the regression results are problematic for model selection. It suggests that even a small amount of measurement noise can lead to significantly different conclusions about the underlying model. In what follows, we quantify this variability while also considering various regression procedures for solving the over-determined linear system Ax = b. Highlighted here are five standard methods: least-square regression (pinv), the backslash operator ($\backslash$), LASSO (least absolute shrinkage and selection operator) (lasso), robust fit (robustfit), and ridge regression (ridge).

这是预期的，因为原始模型是一个带有少量噪声的二次函数。值得注意的是，当添加更多的多项式项时，回归过程中的整体误差实际上增加了，如图(c)所示。因此，仅仅添加更多的项并不能改善误差，这在一开始是反直觉( counter-intuitive)的。


\section{数据驱动的动力系统}

动力系统提供了一个数学框架来描述我们周围的世界，建模了随着时间共同演化的量之间的复杂相互作用(the rich interactions)。正式来说，动力系统关注(concerns)的是分析、预测和理解描述系统状态演化的{\color{red}微分方程系统}或{\color{blue}迭代映射系统}的行为。(Formally, dynamical systems concerns the analysis, prediction, and understanding of the behavior of systems of differential equations or iterative mappings that describe the evolution of the state of a
system.)这种形式足够普遍，能够涵盖各种现象，包括经典机械系统(classical mechanical systems)、电路、湍流、气候科学、金融、生态、社会系统、神经科学、流行病学以及几乎所有随时间演化的系统(nearly every other system that evolves in time)


现代动力系统学的起源可以追溯到庞加莱对行星混沌运动的开创性研究。其根源在于经典力学，并可以看作是从牛顿和莱布尼茨开始的几百年数学建模工作的顶峰。动力系统的历史极其丰富，吸引了几个世纪以来最伟大的头脑的关注，并被应用于无数领域和挑战性问题中。动力系统学作为数学领域中一个最为完整且联系广泛的学科之一，连接了线性代数、微分方程、拓扑学、数值分析和几何学等多个主题。


如今，动力系统学已成为工程学、物理学和生命科学几乎所有领域中建模和分析系统的核心工具。(Dynamical systems
has become central in the modeling and analysis of systems in nearly every field of the
engineering, physical, and life sciences.)


现代动力系统正在经历一场复兴，\textbf{分析推导和基于第一性原理的模型正在逐渐被数据驱动的方法所取代 (with analytical derivations and first principles models giving way to data-driven approaches.)} 大数据与机器学习的结合正在推动科学和工程领域动力系统分析与理解的范式转变。
数据丰富，而物理定律或控制方程仍然难以捉摸，这在气候科学、金融、流行病学和神经科学等问题中尤为显著。即使在诸如光学和湍流等存在控制方程(governing equations)的传统领域，研究人员也越来越多地依赖于数据驱动的分析。许多关键的数据驱动问题，如预测气候变化、通过神经记录理解认知、预测和抑制疾病传播，或控制湍流以实现高效的能源生产和运输，正处于利用数据驱动的动力学发现进展的最佳时机。


特别是，我们将重点关注从数据中发现动力学并寻找使非线性系统适合线性分析的数据驱动表示的关键挑战。(In particular, we will focus on the key challenges of discovering dynamics from data and finding data-driven representations that {\color{red}make nonlinear systems \textbf{amenable} to linear analysis.})


\begin{remark}

The orbits of the planets are open to long-term variations. Modeling the Solar System is a case of the n-body problem of physics, which is generally unsolvable except by numerical simulation. Because of the chaotic behavior embedded in the mathematics, long-term predictions can only be statistical, rather than certain.

The planets' orbits are chaotic over longer time scales, in such a way that the whole Solar System possesses a Lyapunov time in the range of 2-230 million years. In all cases, this means that the positions of individual planets along their orbits ultimately become impossible to predict with any certainty. In some cases, the orbits themselves may change dramatically. Such chaos manifests most strongly as changes in eccentricity, with some planets' orbits becoming significantly more – or less – elliptical.

In calculation, the unknowns include asteroids, the solar quadrupole moment, mass loss from the Sun through radiation and the solar wind, drag of the solar wind on planetary magnetospheres, galactic tidal forces, and effects from passing stars.

\end{remark}

在现代动力系统中，我们通常使用动力系统来模拟现实世界现象，因此与动力系统分析相关的一些高优先级目标包括：

1. **未来状态预测**：在许多情况下，如气象学和气候学中，我们希望预测系统的未来状态。长时间的预测可能仍然具有挑战性。
2. **设计与优化**：我们可能希望调整系统的参数以提高性能或稳定性，例如通过在火箭上放置鳍片来实现。
3. **估计与控制**：通常可以通过反馈主动控制动力系统，利用系统的测量结果来调整控制行为。在这种情况下，通常需要从有限的测量数据中估计系统的完整状态。
4. **可解释性与物理理解**：动力系统的一个更为基础的目标可能是通过分析轨迹和运动方程的解来提供对系统行为的物理洞察和解释。

现实世界的系统通常是非线性的，并在空间和时间上表现出多尺度行为。此外，还必须假设运动方程、参数的规定以及系统测量中存在不确定性。一些系统对这种不确定性更为敏感，因此必须使用概率方法。越来越多的时候，运动的基本方程没有被指定，或者从基本原理出发推导出这些方程可能是不可行的。

本章将讨论识别和分析动力系统的最新数据驱动技术。本章的大部分内容探讨了现代动力系统的两个主要挑战：

1. **非线性**：非线性仍然是分析和控制动力系统的主要挑战，导致复杂的全局动力学。我们之前提到，线性系统可以通过矩阵 \(A\) 的谱分解（即特征值和特征向量）完全表征，从而制定一般的预测、估计和控制程序。对于非线性系统，没有这样的通用框架，而开发这一通用框架是21世纪的一项重大数学挑战。

   非线性动力系统的主要视角考虑了围绕固定点和周期轨道的局部线性化子空间的几何学，连接这些结构的全局异宿轨道和同宿轨道，以及更一般的吸引子。这种几何理论始于庞加莱，已经改变了我们如何模拟复杂系统的方式，其成功在很大程度上归因于一些理论成果，如哈特曼-格罗布曼定理，它们确定了何时以及何处可以用线性动力学来近似非线性系统。因此，在固定点或周期轨道的小邻域内，通常可以应用丰富的线性分析技术。然而，尽管几何视角提供了定量的局部线性模型，但全局分析在固定点和周期轨道之外大多仍是定性和计算性的，这限制了非线性系统的预测、估计和控制理论的发展。

2. **未知动力学**：另一个更为中心的挑战是对于许多现代感兴趣的系统缺乏已知的控制方程。越来越多的研究人员正处理更复杂和更现实的系统，如神经科学、流行病学和生态学中的系统。在这些领域中，缺乏已知的物理定律提供了从中推导运动方程的第一性原理。即使在我们知道控制方程的系统中，如湍流、蛋白质折叠和燃烧，我们仍然难以在这些高维系统中找到模式，从而揭示主导行为沿着哪些内在坐标和粗粒化变量演化。

传统上，物理系统是通过进行理想近似来分析的，然后通过牛顿第二定律推导出简单的微分方程模型。通过利用对称性和巧妙的坐标系，通常可以进行显著简化，正如拉格朗日和哈密顿动力学的成功所证明的那样。随着系统复杂性的增加，这种经典方法正逐渐转变为数据驱动的方法来发现控制方程。

所有模型都是近似的，并且随着复杂性的增加，这些近似往往变得不可靠。确定正确的模型变得更加主观，因此越来越需要自动化的模型发现技术来揭示潜在的物理机制。通常还有一些与动力学相关的潜在变量，但可能未被测量。揭示这些隐藏效应是数据驱动方法的主要挑战。

确定数据中的未知动力学并学习能够线性表示非线性系统的内在坐标是现代动力系统的两个最紧迫的目标。克服未知动力学和非线性的挑战，有望彻底改变我们对复杂系统的理解，并为几乎所有科学和工程领域带来巨大的潜在利益。

本章将在进一步探讨这些问题的过程中，描述一系列应对这些挑战的新兴技术。特别是，有两个关键方法正在定义现代数据驱动的动力系统：

1. **算子理论表示**：为了解决非线性问题，动力系统的算子理论方法正在越来越多地被使用。我们将展示，可以通过无限维但线性的算子来表示非线性动力系统，例如推动测量函数的库普曼算子，以及通过动力学推进概率密度和集合的佩龙-弗罗贝尼乌斯算子。
2. **数据驱动的回归与机器学习**：随着数据越来越丰富，我们继续研究不适合第一性原理分析的系统，回归和机器学习正成为从数据中发现动力系统的重要工具。这是本章中描述的许多技术的基础，包括第7.2节中的动态模态分解（DMD）、第7.3节中的非线性动力学的稀疏识别（SINDy）、第7.5节中的数据驱动的库普曼方法，以及使用遗传编程来从数据中识别动力学。

需要注意的是，本章描述的许多方法和观点是相互关联的，加强和揭示这些关系是正在进行的研究课题。还值得一提的是，另一个主要挑战是许多现代动力系统所涉及的高维性，如人口动力学、大脑模拟和高保真数值离散化的偏微分方程。高维性将在随后的章节中关于降阶模型（ROMs）中得到广泛讨论。


%-----------------------------------------------------------------------

\pdfbookmark[1]{\refname}{ref}

\begin{thebibliography}{}

\bibitem{weinan2021applied} E.~Weinan, Tiejun~Li, and Eric~Vanden-Eijnden. (2021). \emph{Applied stochastic analysis} (Vol. 199). American Mathematical Soc.

\end{thebibliography}


\end{document}